% Original: PDF strana 18, donji slajd.
\slajdid{02-s036}
\begin{frame}[label=02-s036]{Karnoova karta za četiri promenljive}
\setlength{\parskip}{0pt}
\begin{columns}[c,onlytextwidth]
\begin{column}{.48\textwidth}\centering
% element: 02-s036:e01
\begingroup
\def\DEskala{.74}\def\DEpomak{0}
% Zajednički raspored DC/BA. Parametri DEskala i DEpomak dolaze iz omotača.
% Nakon skaliranja oznake imaju 9 pt, indeksi 7 pt (odobren stil s034).
\pgfmathsetmacro{\DEfont}{9/\DEskala}
\pgfmathsetmacro{\DEindeksfont}{7/\DEskala}
\fontsize{\DEfont}{\DEfont}\selectfont
\renewcommand{\small}{\fontsize{\DEfont}{\DEfont}\selectfont}
\scalebox{\DEskala}{%
\begin{karnaugh-map}(label=middle)[4][4][1][$A$][$B$][$C$][$D$]
\foreach \x/\y/\base in {1/3/0,2/3/1,3/3/3,4/3/2,1/2/4,2/2/5,3/2/7,4/2/6,1/1/12,2/1/13,3/1/15,4/1/14,1/0/8,2/0/9,3/0/11,4/0/10}{
 \pgfmathtruncatemacro{\indeks}{\base+\DEpomak}
 \node[anchor=south east,inner sep=1pt,font=\fontsize{\DEindeksfont}{\DEindeksfont}\selectfont] at (\x,\y) {\indeks};
}
\end{karnaugh-map}}

\endgroup

\end{column}
\begin{column}{.48\textwidth}
% element: 02-s036:e02
Po definiciji jediničnog Hemingovog rastojanja, susedna polja razlikuju se u tačno jednom bitu. To važi i preko naspramnih ivica karte.
\par\vspace{3mm}
\Oznaka{Naspramne ivice su susedne}\par\vspace{1mm}
Vertikalno: \(0\leftrightarrow2\), \(4\leftrightarrow6\), \(12\leftrightarrow14\), \(8\leftrightarrow10\).
\par\vspace{2mm}
Horizontalno: \(0\leftrightarrow8\), \(1\leftrightarrow9\), \(3\leftrightarrow11\), \(2\leftrightarrow10\).
\end{column}
\end{columns}
\end{frame}
